% !TeX root = se100_the_ontological_neutrality_theorem.tex
\section{Formal Requirements for Ontological Neutrality}
\label{sec:requirements}

This section formalizes the neutrality requirements
introduced in Section~\ref{sec:introduction}.
All formal results in subsequent sections are conditional
on the definitions, assumptions, and scope conditions stated here.
Unless otherwise noted, all entailment ($\vdash$) is with respect to
classical first-order logic with equality.

The conceptual architecture is summarized in Figure~\ref{fig:architecture}.

% better figure 2026-06-11
\begin{figure}[!ht]
\centering
\begin{tikzpicture}[
    box/.style={rectangle, draw=black, thick, fill=gray!5, minimum width=6.5cm, minimum height=1.2cm, align=center, font=\small},
    arrow/.style={-{Stealth[scale=1.2]}, thick}
]
    % Nodes
    \node[box] (frameworks) {\textbf{Interpretive Frameworks} \\ (Causal and normative reasoning; competing conclusions)};
    \node[box, below=1.5cm of frameworks] (substrate) at (0,-2.5) {\textbf{Neutral Substrate} \\ Framework-invariant referential structure: \\ individuation, co-reference, and persistence for \\ entities, occurrences, and institutional artifacts};

    % Edge
    \draw[arrow] (substrate) -- (frameworks) node[midway, right, font=\scriptsize, xshift=3pt] {extends, but does not revise substrate commitments};
\end{tikzpicture}
\caption{Conceptual architecture of a neutral substrate and interpretive frameworks.}
\label{fig:architecture}
\end{figure}

The interface between layers consists of shared reference to
entities, occurrences, and institutional artifacts,
together with the referential regimes by which they are individuated,
co-referred, and tracked across time.
Interpretive frameworks extend this structure
without revising substrate-layer commitments.

The notion of neutrality used in this paper
is deliberately narrow and operational.
It does not refer to moral neutrality, political impartiality,
epistemic skepticism, or the absence of content.
Rather, neutrality is defined relative to the functional role an ontology
plays when it is used
as a shared substrate across divergent interpretive frameworks.

An ontological substrate is neutral, in the sense relevant here,
if it satisfies two requirements:
interpretive non-commitment and extension stability.
These requirements are independent but jointly necessary.
Together, they characterize the minimal conditions
under which a shared ontology can support accountability
without foreclosing disagreement.

Throughout this section, all requirements are understood
as analytic consequences of this
role-relative definition of neutrality,
not as general design prescriptions for ontologies.

Interpretive non-commitment constrains what the substrate asserts
at its chosen Level of Abstraction.
Extension stability is a global consistency condition:
even a substrate that is silent on contested predicates
can fail neutrality if it encodes background axioms
that conflict with some admissible interpretive framework.

%-------------------------------------------------------
\subsection{Substrate-Layer Commitment}
\label{subsec:substrate-layer-commitment}
%-------------------------------------------------------

\begin{definition}[Substrate-Layer Commitment]
\label{def:substrate-layer-commitment}
A substrate $\mathcal{S}$ \emph{commits to} a proposition $p$
when $\mathcal{S} \vdash p$;
that is, when $p$ is asserted by the substrate
as holding independently of any particular interpretive framework.
A \emph{causal or normative commitment} is a substrate-layer
commitment whose content is a causal relation or normative proposition.
Such commitments do not include the representation of causal or
normative propositions as attributed, provenance-bearing,
reified assertions.
\end{definition}

%-------------------------------------------------------
\subsection{Referential Regimes}
\label{subsec:referential-regimes}
%-------------------------------------------------------

A neutral substrate is not empty.
It must provide stable reference to the entities, occurrences,
and institutional artifacts about which interpretive frameworks reason.
This requires the substrate to fix a referential structure
that remains invariant
across the admissible frameworks under consideration.

\begin{definition}[Referential Regime]
\label{def:referential-regime}
A \emph{referential regime} is the set of conditions under which
a substrate individuates an entity or occurrence,
determines when two references co-refer to the same entity or occurrence,
and tracks that entity or occurrence across time, records,
jurisdictions, or interpretive contexts.
A referential regime therefore comprises three components:
\emph{individuation}, \emph{co-reference}, and \emph{persistence}.
\end{definition}

Referential regimes provide the substrate-level structure
over which interpretive frameworks operate.
They specify what counts as one entity or occurrence,
when different descriptions or records refer to the same one,
and under what conditions that referent is treated as persisting
across time or context.

This terminology separates the referential role of the substrate
from the interpretive role of higher layers.
The substrate may fix a referential regime for an occurrence
without asserting why the occurrence happened,
whether it caused another occurrence,
whether an actor was responsible for it,
or whether it constituted a violation.
Those latter claims belong to interpretive frameworks.

The neutrality theorem therefore presupposes
\emph{referential invariance}:
the relevant referential regimes must be fixable at the substrate layer
and must not themselves be contested
across the admissible frameworks under consideration.
Where individuation, co-reference, or persistence is itself contested,
neutrality in the sense defined here is unavailable until that contestation
is resolved or scoped out.

%-------------------------------------------------------
\subsection{Admissible Frameworks}
\label{subsec:admissible-frameworks}
%-------------------------------------------------------

Admissibility is determined by the domain of application,
not by the substrate.
A framework is admissible if recognized stakeholders
in the domain accept it as a legitimate interpretive stance,
for example, a jurisdiction's legal code,
a scientific community's causal methodology,
or an institution's normative charter.
The substrate must accommodate all such frameworks;
it does not adjudicate among them.

\begin{definition}[Admissible Framework]
\label{def:admissible-framework}
A framework $\mathcal{F}$ is \emph{admissible}
if it is internally consistent
($\mathcal{F} \not\vdash \bot$)
and represents a legitimate interpretive stance
within the domain of application.
The set of all admissible frameworks is denoted $\mathbb{F}$.
Admissible frameworks need not be mutually compatible.
\end{definition}

The theory treats admissibility as a domain-governance parameter,
not as a property determined internally by the substrate ontology.
The substrate is not a court, regulator, or scientific authority.
It does not decide which framework is correct.
It provides a shared referential base over which admissible frameworks
may reason, including when those frameworks reach incompatible conclusions.

%-------------------------------------------------------
\subsection{Framework-Variant Propositions}
\label{subsec:framework-variant}
%-------------------------------------------------------

The first neutrality requirement concerns propositions
whose truth depends on the interpretive framework.

\begin{definition}[Framework-Variant Proposition]
\label{def:framework-variant}
Let $\mathcal{S}$ denote the substrate ontology whose neutrality
properties are under analysis.
A proposition $p$ is \emph{framework-variant}
if there exist admissible frameworks
$\mathcal{F}_1, \mathcal{F}_2 \in \mathbb{F}$
such that
$\mathcal{S} \cup \mathcal{F}_1 \vdash p$
and
$\mathcal{S} \cup \mathcal{F}_2 \vdash \neg p$.
\end{definition}

Framework-variant propositions include, in the domains considered here,
causal and normative conclusions whose truth conditions depend on
legal standards, causal models, institutional rules,
or other interpretive commitments.
Examples include claims that one event caused another,
that an actor was responsible for an occurrence,
or that an action was permitted, prohibited, obligatory,
compliant, noncompliant, justified, or violative.

The governing distinction is between asserting a proposition
as a substrate-layer fact and representing that a framework
has made the proposition as a claim.
The exclusion is at the level of assertion, not vocabulary.
Causal and normative predicates may appear as the opaque content
of reified assertions, reports, or provenance-bearing claim objects.
What neutrality excludes is asserting those
predicates of substrate individuals as substrate-layer facts.

\begin{assumption}[Contestability of Causal and Normative Propositions]
\label{asm:contestability}
In the domains to which the neutrality requirements are applied,
causal and normative propositions are contestable across admissible frameworks.
In particular, if a substrate $\mathcal{S}$
asserts a causal or normative proposition $p$
as a substrate-layer commitment,
then there exists an admissible framework $\mathcal{F} \in \mathbb{F}$ such that
$\mathcal{S} \cup \mathcal{F} \vdash \neg p$.
\end{assumption}

This assumption is a scope condition on the theorem.
It does not claim that causation is unreal,
that normative judgment is arbitrary,
or that every domain contains an active dispute about every claim.
It states that, in the class of domains addressed by this paper,
causal and normative conclusions are not guaranteed to remain invariant
across admissible interpretive frameworks.
Accordingly, asserting such conclusions at the substrate layer
risks converting the substrate into one framework's interpretation.

%-------------------------------------------------------
\subsection{Interpretive Non-Commitment}
\label{subsec:interpretive-non-commitment}
%-------------------------------------------------------

The first neutrality requirement is interpretive non-commitment.
A substrate ontology satisfies this requirement
if it does not assert any proposition
whose truth value can vary across admissible interpretive frameworks.

\begin{definition}[Interpretive Non-Commitment]
\label{def:interpretive-non-commitment}
A substrate ontology $\mathcal{S}$ satisfies \emph{interpretive non-commitment}
if it does not assert any framework-variant proposition:
there is no proposition $p$ such that $\mathcal{S} \vdash p$
and $p$ is framework-variant with respect to $\mathbb{F}$.
\end{definition}

Accountability systems routinely involve disagreements
about what occurred, why it occurred, and how it should be evaluated.
Legal frameworks may disagree about whether an action was permitted or prohibited;
analytic frameworks may disagree about whether one event caused another;
political frameworks may disagree about responsibility or intent.
These disagreements are not merely epistemic gaps
to be filled with additional data.
They reflect differences in background assumptions,
governing rules, and interpretive commitments.

A neutral substrate,
under this definition,
must therefore refrain from
settling such questions at the substrate layer.
If a substrate ontology asserts that a particular action was
obligatory, forbidden, permitted, compliant, or violative,
it commits the ontology to a particular normative framework.
Similarly, if it asserts that one event caused another,
it commits the ontology to a particular causal model or explanatory framework.
This does not imply that causation is subjective or unreal.
Rather, causal conclusions depend on explanatory frameworks
such as statistical, mechanistic, or legal standards of causation
that may legitimately differ across admissible frameworks.
In both cases, asserting such conclusions as substrate-layer commitments
violates interpretive non-commitment.

Interpretive non-commitment does not require the substrate ontology
to ignore normative or causal discourse.
Rather, it requires that the substrate's own assertions
remain restricted to framework-invariant referential structure.
Causal and normative claims may be represented as claims,
with attribution and provenance,
but their truth is evaluated only within interpretive frameworks
layered above the substrate.

%-------------------------------------------------------
\subsection{Extension Stability}
\label{subsec:extension-stability}
%-------------------------------------------------------

The second neutrality requirement is extension stability.
A substrate ontology satisfies this requirement if it remains consistent
when extended separately by each admissible interpretive framework.

\begin{definition}[Extension Stability]
  \label{def:extension-stability}
  A substrate ontology $\mathcal{S}$ satisfies \emph{extension stability}
  if for all admissible frameworks $\mathcal{F} \in \mathbb{F}$,
  $\mathcal{S} \cup \mathcal{F} \not\vdash \bot$.
\end{definition}

Extension stability follows from the functional role of a neutral substrate,
which must remain usable by independently developed
interpretive frameworks whose conclusions may diverge
and whose development cannot be centrally coordinated.

Extension stability does not require that all admissible interpretive frameworks
be simultaneously combined into a single globally consistent theory.
Rather, the requirement is \emph{pairwise compatibility with the substrate}:
for each admissible framework $\mathcal{F} \in \mathbb{F}$ considered independently,
the combined theory $\mathcal{S} \cup \mathcal{F}$ must remain consistent.
The substrate need not reconcile frameworks with one another.
It must avoid assertions that require revision when
an admissible framework disagrees.

The substrate is permitted to be extended by additional structures,
records, or interpretive layers, provided such extensions do not retract,
revise, or contradict substrate-layer assertions.
A violation occurs when accommodating an admissible framework
requires modification of the substrate itself.
Disagreement between frameworks, including disagreement about causal,
normative, or explanatory conclusions, is expected and permitted.
What is prohibited is embedding substrate-layer commitments
that make such disagreement impossible to accommodate
without revising the substrate.

Consider a substrate ontology $\mathcal{S}$ used by two interpretive frameworks,
$\mathcal{F}_1$ and $\mathcal{F}_2$, that contradict each other
($\mathcal{F}_1 \cup \mathcal{F}_2 \vdash \bot$).
For $\mathcal{S}$ to remain stable, it must be separately consistent
with each framework:
$\mathcal{S} \cup \mathcal{F}_1 \not\vdash \bot$ and
$\mathcal{S} \cup \mathcal{F}_2 \not\vdash \bot$.
The substrate need not reconcile the frameworks.
It must simply avoid assertions that either framework rejects.

Now suppose $\mathcal{S}$ commits to a causal or normative proposition $p$,
and some admissible framework $\mathcal{F}_2$ rejects it
($\mathcal{S} \cup \mathcal{F}_2 \vdash \neg p$).
Since $\mathcal{S} \vdash p$,
the combined theory entails both $p$ and $\neg p$:
$\mathcal{S} \cup \mathcal{F}_2 \vdash \bot$.
The substrate must then be revised to accommodate
$\mathcal{F}_2$, violating extension stability.

Now suppose $\mathcal{S}$ commits to a causal or normative proposition $p$,
Such a substrate cannot function as a neutral foundation
for the class of frameworks that includes such rejection,
because accommodating the rejecting framework would require
revision of the substrate.

%-------------------------------------------------------
\subsection{Levels of Abstraction and Organization}
\label{subsec:loas-and-loos}
%-------------------------------------------------------

The neutrality requirements constrain both the
Level of Abstraction (LoA) and the
Level of Organization (LoO) of the architecture.

In the Method of Levels of Abstraction,
an LoA is determined by the observables made available
for representation~\citep{floridi2008levels,floridi2011philosophy}.
For a substrate ontology, these observables correspond
to the predicate symbols and assertion forms available
at the foundational layer.
Neutrality requires that the substrate's own assertions
be restricted to framework-invariant referential predicates:
predicates supporting referential regimes,
including individuation, co-reference, persistence,
provenance, temporal ordering, participation,
containment, and other structural relations
that assert neither causal efficacy nor normative force.

This is not a prohibition on mentioning causal or normative content.
Causal and normative predicates may occur within reified assertions
as opaque content, for example
$\mathsf{Asserts}(\mathcal{F}, \varphi)$,
where $\varphi$ is a causal or normative claim.
The substrate may assert that a framework made a claim.
It may not assert the claim's truth as a substrate-layer fact.

The LoO characterizes the architectural hierarchy of the system:
the ordering of layers by direction of dependence.
The neutral substrate occupies the foundational organizational layer.
Higher interpretive layers may extend the substrate
by adding causal models, normative rules, legal standards,
or evaluative conclusions.
They may not mutate, retract, contradict,
or revise substrate-layer commitments.

This directional dependence is essential.
Identity, co-reference, and persistence are established
at the substrate layer before interpretive frameworks are applied.
Explanation, evaluation, and judgment operate on referents
whose referential regimes have already been fixed.
Changes in higher-level interpretive logic must therefore produce
new interpretive assertions or revised framework-level conclusions,
not revisions to the substrate itself.

%-------------------------------------------------------
\subsection{Scope of Neutrality}
\label{subsec:scope}
%-------------------------------------------------------

Neutrality, as defined here, is not global.
A neutral substrate is not neutral with respect to
the referential regimes it asserts.
It must take substantive positions on individuation,
co-reference, and persistence,
because accountability requires stable referents.
These commitments are necessary for accountability
and do not threaten neutrality insofar as they are fixed
as part of the shared substrate scope
and remain invariant
across the admissible frameworks under consideration.

Neutrality is therefore scoped specifically to interpretive commitments:
those aspects of representation whose truth depends on
legal, normative, causal, or analytic frameworks
that may legitimately disagree.
This restriction concerns what the substrate may assert,
not whether causal or normative relations exist in the world.
Such relations may be represented and reasoned about within
interpretive frameworks layered above the substrate.
The exclusion of such commitments from the substrate
is not a matter of preference or conservatism;
it is a functional requirement derived from
the role the ontology is intended to play
as a neutral substrate across divergent interpretive frameworks.

A further scope condition follows.
The neutrality requirements assume that the substrate's
framework-invariant referential structure, including its
referential regime, is itself invariant across admissible frameworks.
If frameworks disagree about
individuation, co-reference, or persistence---for example,
whether two references denote the same actor,
whether an occurrence should be treated as one event or many,
or whether a particular institutional artifact persists
across legal or organizational change---then
the substrate cannot be neutral with respect to those frameworks
until that referential contestation is resolved or scoped out.
This is not a defect of the framework but a boundary condition:
neutrality presupposes a shared ontological ground on which interpretation can vary.

Together, interpretive non-commitment and extension stability
define the functional boundaries of a neutral substrate.
These requirements show that neutrality is not a lack of content,
but a categorical constraint on where interpretive commitments may occur.
The next section establishes that,
under the contestability assumption and referential-invariance
scope condition stated here,
a substrate ontology satisfies neutrality if and only if
its foundational layer is restricted to framework-invariant
referential structure and makes no causal or normative commitments.
